\section{The type theory}
\label{appendix:type-theory}
The symbol $\equiv$ denotes definitional equality in the theory for both types and terms. We make use of the following forms of judgements:
\begin{align*}
    \vdash \Gamma \; \Ctxt, \qquad \Gamma \vdash A \Type, \qquad \Gamma \vdash A \equiv B \Type, \qquad \Gamma \vdash a : A, \qquad \Gamma \vdash a \equiv b : A.
\end{align*}
For brevity, we omit the standard rules stating that definitional equality is a congruence relation with respect to all type- and term formers. We may also assume weakening and substitution rules with respect to all judgements, which are both admissible \cite{Hofmann_1997}.

\input{chapters/syntax}

Function extensionality is encoded as in \cite{hottbook}: given $f, g : \Pi x : A.\,B(x)$, there is a map \[
    (f = g) \to \Pi x : A.\, (f(x) = g(x)),
\] and the function extensionality axiom states that this map is an equivalence.
Similarly, we note that for small types \(A\), \(B(i,x)\), and \(C(i)\) there are canonical maps, for which the parametricity axioms state that each such map is an equivalence: 
\begin{align*}
    %A &\simeq \forall i.\, A &
    %\exists i \. A &\simeq A &
    %i \notin \text{fv}(A) \\
    \Pi x : A.\,\exists i . B(x, i)&\leftarrow\exists i \. \Pi x : A. \, \exists j < i .B(x, j), \\
    \Sigma x:A.\,\forall i . B(x, i) &\to \forall i \. \Sigma x : A.\,B(i, x), \\
    \exists i \.  C(i)&\leftarrow \exists i \. \exists j < i \.C(j), \\
    \forall i \.  C(i) &\to \forall i \. \forall j < i \.C(j).
\end{align*} 

\newpage

\section{Encoding}
\label{appendix:encoding}
\subsection{Inductive types}

\begin{proof}[Proof of \autoref{lemma:inductive-sized-encoding}]
    The propositional \( \beta \)-rule for the fixpoint operator gives us the following equality:
    \begin{align*}
        \mu^\Diamond F \, i &\coloneqq \fix (\lambda i \. \lambda X\. F(\Diamond_i X)) \, i \\
        &= (\lambda i \. \lambda X\. F(\Diamond_i X)) \, i \, (\lambda r < s \. \fix (\dots) \, r) \\
        &\equiv F(\Diamond_s \, \fix (\lambda i \. \lambda X\. F(\Diamond_i X)) ) \\
        & =: F [\Diamond (\mu^\Diamond F)] \, i
    .\end{align*}
    The structure map \( \In^\Diamond : F[\Diamond (\mu^\Diamond F)] \sto \mu^\Diamond F \) is then given by the pointwise transport operation on the path given by $\fix^\beta f \, s$. This map is invertible, and we denote its inverse by $\out^\Diamond$. Given an algebra \( (A, k_A) \), the term \( \fold^\Diamond \, A \, k_A : \mu^\Diamond F \sto A\) is also defined using the fixpoint operator as follows:
    \begin{align*}
        \fold^\Diamond \, A \, k_A  &\coloneqq \fix \, (\lambda i \. \lambda \fold^{\Diamond'} \. \lambda m \. (k_A \circ F[\Diamond \fold^{\Diamond'}] \circ \mathsf{out}^\Diamond) \, i \, m).
    \end{align*}
    Thus, for $i : \Size$ and $m : \mu^\Diamond F \, i$, by $\fix^\beta$ we have
    \[
    (\fold^\Diamond \, A \, k_A) \, i \, m = (k_A \circ F[\Diamond \fold^{\Diamond}] \circ \mathsf{out}^\Diamond) \, i \, m.
    \]
    We obtain a path $\fold^\Diamond \, A \, k_A \circ \In^\Diamond = k_A \circ F[\Diamond (\fold^\Diamond \, A \, k_A)]$ using $\fix^\beta f$ and the canonical proof that $\out^\Diamond$ is the inverse of $\In^\Diamond$. 
    To show that the following is contractible: \[\Sigma{(h : \mu^\Diamond F \to A)} \. (h \circ \In^\Diamond = k_A \circ F[\Diamond h])\] is equivalent to show that the following is contractible: \[\Sigma{(h : \mu^\Diamond F \to A)} \. (h = k_A \circ F[\Diamond h] \circ \out^\Diamond).\] 
    By function extensionality, this is then equivalent to showing that the following type is contractible: \[\Sigma{(h : \mu F \to A)} \, \forall i \. \Pi{(m : \mu^\Diamond F \, i)} \. (h \, i \, m = (k_A \circ Fh \circ \out^\Diamond) \, i \, m).\]
    Since $\fold^\Diamond \, A \, k_A$ is defined via the fixpoint operator, this follows from \autoref{lemma:unique-fix}.
\end{proof}

\subsection{Coinductive types}
\begin{definition}[Coalgebra]
    Given an endofunctor \( F : U \to U \), an \emph{F-coalgebra} is a pair \( (C, k_C) \), consisting of a type \( C : U \) and a structure map \( k_C : C \to FC \). A map of coalgebras \( (h, s_h) : (C, k_C) \to (D, k_D) \) is a pair of a map \( h : C \to D \) and a path \( s_h : k_D \circ h = Fh \circ k_C\). We write $\FCoalg$ for the type of $F$-algebras.

A \emph{terminal $F$-coalgebra} is an coalgebra \( (\nu F, \out) \), such that for any coalgebra \( (C, k_C) \), the following type is contractible: 
\[\Sigma{(h : C \to \nu F)} \, (\out \circ h = Fh \circ k_C).\]
\end{definition}
We define \emph{size-indexed \( F \)-coalgebras} as coalgebras for \( F[\Box -] \). Explicitly, a size-indexed \( F \)-coalgebra is a pair \( (C, k_C) \) where \( C : \Size \to U \) and \( k_C : C \sto F[\Box C] \).
%We will use the notation \( F[\Box{-}] \)-coalgebra to mean \( F[\Box -] \)-coalgebra where convenient, and write $\sFCoalg$ for the type of $F[\Box{-}]$-coalgebras. 

To encode size-indexed coinductive types, we construct terminal coalgebras for size-indexed functors. This is done in the following lemma, where we make use of the fixpoint operator:
\begin{lemma}\label{lemma:coinductive-sized-encoding}
    Let \( F \) be a functor on small types. The size-indexed functor \( F[\Box -] \) has a terminal coalgebra \( (\nu^\Box F, \out^\Box) \), where the first component is given by:
\[
    \nu^\Box F \, i \coloneqq  \fix (\lambda i. \lambda  X. F(\Box_i X)) \, i.
\] 
\end{lemma}
\begin{proof}[Proof Sketch]
    We note that the propositional \( \beta \)-rule for the fixpoint operator gives us the following equality:
    \begin{align*}
        \nu^\Box F \, i &\coloneqq \fix (\lambda i \. \lambda X\. F(\Box_i X)) \, i \\
        &= (\lambda i \. \lambda X\. F(\Box_i X)) \, i \, (\lambda r < s \. \fix (\dots) \, r) \\
        &\equiv F(\Box_i \, \fix (\lambda s \. \lambda X\. F(\Box_s X)) ) \\
        & =: F [\Box (\nu^\Box F)] \, i
    .\end{align*}
    The structure map \( \out : \nu^\Box F \sto F[\Box (\nu^\Box F)] \) is then given by the pointwise transport operation. Given a coalgebra \( (C, k_C) \), the term \( \unfold^\Box \, C \, k_C : C \sto \nu^\Box F \) is also defined using the fixpoint operator in the same way as in the algebra case, and can be shown to be an coalgebra morphism. The fact that it is a terminal coalgebra then follows from the contractibility of \( \fix \!\) described in \autoref{lemma:unique-fix}, together with function extensionality.
\end{proof}

As for algebras, we can lift a coalgebra to a size-indexed coalgebra. We reuse the notation $T$ from algebras.
\begin{lemma}
    Every \( F \)-coalgebra \( (C, k_C) \) gives rise to an \( F[\Box{-}] \)-coalgebra \( (TC, k_{TC}) \), where the first component is defined by 
    \[
        TC \coloneqq \lambda i. C : \Size \to U.
    \] 
This gives rise to a functorial map \( T : \FCoalg \to \sFCoalg \). 
\end{lemma}
\begin{proof}[Proof Sketch]
    We require a structure map \( k_{TC} : TC \sto F[\Box TC]\). Noting that \( i \notin \text{fv}(C) \), we can define the following auxiliary term:
    \begin{align*}
        \mathsf{insert}& : \Pi{C:\U} \. TC \sto \Box TC, \\
        \mathsf{insert}&_C \, i \, c \coloneq  \lambda j < i \. c
    \end{align*}
    This allows us to define $k_{TC} \, i \coloneq  F(\mathsf{insert}_C \, i) \circ k_C$, which is well-typed.
    Given a coalgebra morphism $f : (C, k_C) \to (D, k_D)$, we define $Tf : ({TC}, k_{TC}) \to ({TD}, k_{TD})$ by
    \begin{align*}
        Tf &: TC \sto TD, \\
        Tf & \coloneq  \lambda i . f.
    \end{align*}

    By unfolding the definitions, it is clear that $\mathsf{insert}$ is natural in $C$. As in the algebra case, this is sufficient to show that $Tf$ satisfies the required equality for a coalgebra morphism. Thus, $T$ as defined here is a functorial map from $F\text-\sf{CoAlg}$ to $F[\Box{-}]\text-\sf{CoAlg}$.
\end{proof}

The analogous property that we require a functor to satisfy in order to construct a terminal coalgebra is given as follows:

\begin{definition}
\label{weakly-commute-forall}
    Let $\F$ be a functor on the category of small types. We say that $\F$ \emph{weakly commutes} with $\forall$ if the following canonical map
    \begin{align*}
    \mathsf{can}_C &: \F (\forall i .C \, i) \to (\forall i \. F [\Box C] \, i), \\
    \mathsf{can}_C &\, c' \, i \coloneq  F \, (\psi_C \, i) \, c' .
    \end{align*}
    is an equivalence for each $C : \Size \to \U$, where we define $\psi_C$ as:
    \begin{align*}
        \psi_C &: \forall i \. (\forall j \. C \, j \to \forall j < i \. C \, j), \\
        \psi_C & \, i \, c' \, j \coloneq  c' j.
    \end{align*}
\end{definition}

For functors that satisfy this property, we can define a functor $\forall : F[\Box{-}]\text-\sf{CoAlg} \to F\text-\sf{CoAlg}$ extending the action of the universal quantifier on size-indexed types as follows.
\begin{proposition}
    Given a functor $F$ on the category of small types that weakly commutes with $\forall$, every $F[\Box-]$-coalgebra $(C, k_C)$ gives rise to a type $\forall i . A \, i$ which is the first component of an $F$-coalgebra. Moreover, this operation is functorial.
\end{proposition}
\begin{proof}
    First, we note that the universal quantifier works on maps too, meaning that for any map $f : C \sto D$, we define a map betewen universally quantified types as:
    \begin{align*}
        &\forall i . C \, i \to \forall i . D \, i, \\
        &\forall f \coloneq  \lambda i . f \, i \, (g \, i).
    \end{align*}

    Now let $\mathsf{can}_C^{-1} : \forall i \. F [\Box C] \, i \to \F (\forall i .C \, i) $ be the isomorphism obtained from $F$ weakly commuting with $\forall$.
    For an $F[\Box-]$-coalgebra $(C, k_C)$, where $k_C : C \sto F[\Box C]$, we obtain the $F$-coalgebra $(\forall i . \, C \, i \, , k_{\forall C})$ where we define
    \begin{align*}
        k_{\forall C}& : \forall i . \, C \, i \to F(\forall i . \, C \, i), \\
        k_{\forall C}& \coloneq  \mathsf{can}_C^{-1} \circ \forall  k_C.
    \end{align*}

    Given an coalgebra morphism $f: (C, k_C) \to (D, k_D)$, we obtain the map
    $\forall f : \forall i . C \, i \to \forall i . D \, i$. The fact that this constitutes a coalgebra morphism follows from the fact that $\mathsf{can}_C$ is natural in $C$.
\end{proof}

Dually to the case for initial algebras, we are able to prove that for functors that weakly commute with $\forall$, their terminal coalgebra is obtained by universal quantification over the terminal size-indexed coalgebra.
\begin{theorem}
    Let $\F$ be a functor on the category of small types that weakly commutes with the universal quantifier. Then $\F$ has a terminal coalgebra with first component:
    \[
    \nu  \F \coloneq  \forall i . (\nu^\Box F) \, i.
    \]
\end{theorem}
\begin{proof}
This argument is also very similar to the algebra case.
The categorical intuition is that this follows from the fact that the functor $\forall : \sFCoalg \to \FCoalg$ of the previous proposition is the right adjoint of $T : \FCoalg \to \sFCoalg$.
Since right adjoints preserve terminal objects, we obtain the fact that $\forall i. (\nu^\Box F) \, i$ is the terminal coalgebra of $F$, with the structure map $\out \coloneq  \mathsf{can}^{-1}_{\nu^\Box F} \circ \forall \, \In^\Box$, which is also an equivalence.

Working this out in more detail: given a type $C : U$ and a size-indexed type $D : \Size \to U$, there is an isomorphism between maps of type $TC \sto D$ and those of type $C \to \forall i . D \, i $, which is given by swapping arguments. This extends to an isomorphism between size-indexed coalgebra maps of type $(TC, k_{TC}) \to (D, k_{D})$ and coalgebra maps of type $(C, k_C) \to (\forall i . D \, i, k_{\forall D})$. Thus, defining a coalgebra map $(C, k_C) \to (\nu \F, \out) $ is equivalent to defining an size-indexed coalgebra map $(C, k_{TC}) \to (\nu^\Box F, \out^\Box)$. Such a coalgebra map uniquely exists by the fact that $(\nu^\Box F, \out^\Box)$ is the terminal size-indexed coalgebra.  
\end{proof}

Finally, we can show that all polynomial functors satisfy the requirement of weakly commuting with the universal quantifier. As before, the equivalences entail the following chain of equivalences, which corresponds precisely to the canonical map denoted in \autoref{weakly-commute-forall}, noting that $i \notin \text{fv}(A)$ and $i \notin \text{fv}(B(a))$ for $a : A$:
\begin{align*}
    P_{A,B} (\forall i .X \, i) &\coloneq  \Sigma {a : A} \. (B (a) \to \forall i .X \, i) \\
    &\simeq \Sigma {a : A} \. (B (a) \to \forall i . \forall j < i . X \, j) \\
    &\simeq \Sigma {a : A} \. \forall i \. B(a) \to \forall j < i . X \, j \\
    &\simeq \forall i \. \Sigma {a : A} \. B(a) \to \forall j < i . X \, j
    =: (\forall i \. P_{A,B} [\Box X] \, i).
\end{align*}
Thus, for all polynomial functors, the terminal coalgebra is constructed via universal quantification.


\section{Model of the theory}
\label{appendix:model}
In this section we give more details on how the category of assemblies fits in the CwF framework, and how our additional type formers are interpreted in the model. Recall that we work with the pca of partial recursive functions. We also assume a computable bijective pairing function $\langle \_\, , \_ \rangle : \bN^2 \to \bN$ is given, together with projections $\pr_1, \pr_2 : \bN \to \bN$. An enumeration of partial recursive functions and an appropriate pairing function can be found in many treatments of recursion theory. Another important feature of a pca is that it has a notion of abstraction similar to $\lambda$-abstraction in the lambda calculus. For a variable $x$ and a term $t$, there exists another term $\Lambda x . t$ defined by recursion on terms, which behaves exactly as $\lambda$-abstraction. See, for example, \cite{oosten} for more details on partial combinatory algebra.

\subsection{Assemblies as a CwF}
 The CwF framework can be defined in different ways, we make use of the definitions from Hofmann \cite{Hofmann_1997}, in which the construction of a CwF begins with the following structure:
\begin{itemize}
	\item A category $\C$ of semantic contexts and context morphisms;
	\item For each $\Gamma \in \C$, a collection $\Ty(\Gamma)$ of semantic types in context $\Gamma$.
	\item For each $\Gamma \in \C$ and $A \in \Ty(\Gamma)$, a collection $\Tm(\Gamma, A)$ of semantic terms of type $A$ in context $\Gamma$.
\end{itemize}
We will often leave out ``semantic'' where this does not cause confusion. The category $\Asm$ is the category of contexts. Given an assembly $\Gamma$, a type $A \in \Ty(\Gamma)$ will be a family of assemblies over $\Gamma$ We will often write such a family as $A = (A_\gamma)_{\gamma \in \Gamma}$. Finally, a term $a \in \Tm(\Gamma, A)$ is a dependent morphism $a : (\gamma \in \Gamma) \to A_\gamma$, often written $a = (a_\gamma)_{\gamma \in \Gamma}$, which means that $a$ is tracked by some $e \in \bN$, i.e.
\[
\forall n \in \bN . \forall \gamma \in \Gamma \. n \Vdash_\Gamma \gamma \implies (e \cdot n) \Vdash_{A_\gamma} a_\gamma.
\]

A CwF requires a notion of \emph{semantic substitution}, given by an operation on types and one on terms. Given a context morphism $f :  \Delta \to \Gamma$, there should be an operation $-[f] : \Ty(\Gamma) \to \Ty(\Delta)$, and for each $A \in \Ty(\Gamma)$ an operation $-[f] : \Tm(\Gamma, A) \to \Tm(\Delta, A[f])$. Moreover, these operations should be compatible with identities and composition \cite{Hofmann_1997}. 
For assemblies, both operations are defined by precomposition with $f$. Given an assembly morphism $f :  \Delta \to \Gamma$ and a type $A \in \Ty(\Delta)$, we define $A[f]_\delta = A_{f(\delta)}$. Similarly, if $a \in \Tm(\Gamma, A)$, we define $(a[f])_\delta = a_{f(\delta)}$.

A CwF should be able to interpret the \emph{context formation} rules of the theory. The empty context $\diamond$ is modelled by the terminal object in $\C$. In the category of assemblies, the terminal object is given by $\1 = (\{*\}, \nabla \{*\})$. To interpret context extension, we require for each $\Gamma \in \C$ and $A \in \Ty(\Gamma)$ an extended context $\Gamma.A \in \C$, together with context morphisms $\p_A : \Gamma.A \to \Gamma$ and $\v_A \in \Tm(\Gamma.A, A[\p_A])$, corresponding to first and second projections respectively. Context extension is defined as follows for assemblies:
\begin{definition}[Context extension]
	Given an assembly $\Gamma$ and $A \in \Ty(\Gamma)$, we define the context extension $\Gamma.A$ as:
	\begin{align*}
		& \Gamma.A = \{ (\gamma, a) \mid \gamma \in \Gamma \land a \in A_\gamma\}, \\
		& \langle n_1, n_2 \rangle \Vdash_{\Gamma.A} (\gamma, a) \Iff n_1\Vdash_\Gamma \gamma \and n_2 \Vdash_{A_\gamma} a.
	\end{align*}
	The morphisms $\p_A : \sigma(\Gamma, A) \to \Gamma$ and $\v_A \in \Tm(\Gamma.A, A[\p_A])$ are given by $(\gamma, a) \mapsto \gamma$ and $(\gamma, a) \mapsto a$ respectively, which are assembly morphisms tracked by the first and second projections on the realisers. It is moreover clear that in this model, $A[\p_A]_{(\gamma, a)} = A_\gamma$. 
\end{definition}

Finally, a CwF requires an operation to extend a context morphism by a term. That is, given a context morphism $f : \Delta \to \Gamma$, a type $A \in \Ty(\Gamma)$, and a term $a \in \Tm(\Delta, A[f])$, there should be a morphism $\langle f, a \rangle : \Delta \to \Gamma.A$ satisfying several equations, as found in \cite{Hofmann_1997}. In the realisability model, this is defined as $\langle f, a \rangle (\delta) = (f_\delta, a_\delta)$. Note that if $d$ tracks $f$ and $e$ tracks $M$, then $\langle d, e \rangle$ tracks $\langle f, a \rangle$, so this is indeed an assembly morphism. If $a \in \Tm(\Gamma, A)$, then we trivially have that $a \in \Tm(\Gamma, A[\id_\Gamma])$, so we get a context morphism $\langle \id_\Gamma, a \rangle : \Gamma \to \Gamma . A$, where $\langle \id_\Gamma, a \rangle (\gamma) = (\gamma, a_\gamma)$. We will denote this morphism by $\overline{a} : \Gamma \to \Gamma.A$, as it will be useful later.

This completes the construction of a category with families from the category of assemblies, allowing us to interpret the structural rules of our dependent type theory in assemblies. What remains is to give interpretations for the different types and type formers.

\subsection{Universe type}
The universe type is interpreted by giving the set of all PERs the complete realisability relation.
\begin{definition}[Universe type]\label{def:universe}
	For any context $\Gamma$, there is a type $\U \in \Ty(\Gamma)$ defined as $U_\gamma = \nabla \Per$, viewing $\Per$ as the set of all PERs. 
    Moreover, there exists a type $\El \in \Ty(\Gamma.\U)$ defined by $\El_{(\gamma, R)} = \iota(R)$, where $\iota(R)$ is the PER $R$ viewed as an assembly. Given a type $A \in \Tm(\Gamma, \U)$ in the universe, we obtain a type $\El[\overline{A}] \in \Ty(\Gamma)$, which we will denote by $\El A$.
\end{definition}

\subsection{Empty, unit and boolean Types}
The empty-, unit-, and boolean types are non-dependent types, and thus are are interpreted as constant families. We write $\0$ and $\1$ for the initial and terminal PERs respectively. Note that for these types, we implicitly assume an induction principle, which is also not explicitly specified in the type theory, but is easily defined in the interpretation.
\begin{definition}[Empty type]
The code for the empty type $\bot^\U \in \Tm(\Gamma, \U)$ will be interpreted as the initial PER for each $\gamma \in \Gamma$:
\[
\bot^\U(\gamma) = \0.
\]
This definition allows us to semantically define the empty type as $\bot \coloneq  \El (\bot^\U)$, where $\El$ denotes the semantic decoder given in \autoref{def:universe}.
\end{definition}

\begin{definition}[Unit type]
The small unit type $\top^\U \in \Tm(\Gamma, \U)$ will be interpreted as the terminal PER for each $\gamma \in \Gamma$:
\[
\top^\U(\gamma) = \1.
\]
As with the empty type, the actual unit type is interpreted as $\El (\top^\U)$, i.e. the terminal assembly for each $\gamma$, making the required type equality true in the interpretation.
The unique term $\star \in \Tm(\Gamma, \top)$ is then given by the unique morphism from $\Gamma$ to the terminal assembly.
\end{definition}

\begin{definition}[Boolean type]
We define $\mathsf{Bool}^\U$ as the following PER:
\begin{align*}
    \forall m, n \in \bN \. \langle 0, m \rangle \sim_{\mathsf{Bool}^\U_\gamma} \langle 0, n \rangle, \\
    \forall m, n \in \bN \. \langle 1, m \rangle \sim_{\mathsf{Bool}^\U_\gamma} \langle 1, n \rangle.
\end{align*}
Again, we define the interpretation of the type $\mathsf{Bool}$ as $\El (\mathsf{Bool}^\U)$. The terms $\mathsf{ff, tt} \in \Tm(\Gamma, \mathsf{Bool})$ are defined to be the constant morphisms $\gamma \mapsto \{\langle 0, m \rangle \mid m \in \bN \}$ and $\gamma \mapsto \{\langle 1, m \rangle \mid m \in \bN \}$, realised by $\Lambda k . \langle 0, k \rangle$ and $\Lambda k . \langle 1, k \rangle$ respectively.
\end{definition}

\subsection{Identity type and function extensionality}
\begin{definition}[Identity type]
We interpret the identity type as follows: for any type $A \in \Ty(\Gamma)$ and terms $a, a' \in \Tm(\Gamma, A)$, we have a type $\Id(A, a, a') \in \Ty(\Gamma)$ defined as 
\[
\Id(A, a, a')_\gamma = \begin{cases}
	\1 & \text{ if $a = a'$},\\
	\0 & \text{ if $a \neq a'$}.
\end{cases}
\]
For each term $a \in \Tm(\Gamma, A)$, we have a term $\refl_a \in \Tm(\Gamma, \Id(A, a, a))$ defined, closely following the set-theoretic model, by $\refl_a(\gamma) = *$.
\end{definition}

Note that the realisability model thus validates uniqueness of identity proofs, since the identity types are either empty or a singleton, but we do not add this axiom to the type theory. 
Informally, function extensionality holds in the realisability model since it is set-theoretic, so (dependent) functions are equal when they are pointwise equal.

\subsection{Dependent product}
To interpret both large and small dependent products in assemblies, we are required to define separate large and small type- and term constructors. This is due to the fact that the inclusion of the small dependent product into assemblies is only isomorphic to the large dependent product, not equal.
To satisfy the definitional type equality
\[
\El (\Pi^U x:A.\, B(x)) \equiv \Pi_{x:\El A} (\El B(x)),
\]
a case distinction is the required in the actual interpretation. We only give the definitions for the product type. First we define the large product and its associated terms.

\begin{definition}[Dependent product]
	Given types $A \in \Ty(\Gamma)$ and $B \in \Ty(\Gamma.A)$, the type $\Pi(A, B) \in \Ty(\Gamma)$ consists of all dependent morphisms from $A$ to $B$, defined as follows:
	\begin{align*}
		&\Pi(A, B)_\gamma = \{ f \in \Pi_{a \in A_\gamma}  B_{(\gamma, a)} \mid \exists e \. e \Vdash_{\Pi(A, B)_\gamma} f \},\\
		&e \Vdash_{\Pi(A, B)_\gamma} f \Iff \forall n \in \bN, \forall a \in A_\gamma \. n \Vdash_{A_\gamma} a \implies e \cdot n \Vdash_{B_{(\gamma, x)}} f (a).
	\end{align*}
	Given a term $b \in \Tm(\Gamma.A, B)$, we define a term $\abs^\Pi_{A, B} (b) \in \Tm(\Gamma, \Pi(A, B))$ as
	\[
	\abs^\Pi_{A, B} (b)(\gamma) = \lambda a \in A_\gamma \. b_{(\gamma, a)}.
	\]
	Note that $\abs^\Pi_{A, B}$ is tracked by $\Lambda m . \Lambda z . \Lambda n\. m \cdot \langle z, n \rangle$. 
	Conversely, given a term $f \in \Tm(\Gamma, \Pi(A, B))$ and $a \in \Tm (\Gamma, A)$, we have a term $\app^\Pi_{A, B}(f, a) \in \Tm(\Gamma, B[\overline{a}])$ defined as follows
	\[
	\app^\Pi_{A, B}(f, a)(\gamma)  = f_\gamma(a_\gamma).
	\]
	Note that $\app^\Pi_{A, B}(f, a)$ is tracked by $\Lambda e . \Lambda n . \Lambda z \. (e \cdot z) \cdot (n \cdot z)$.
\end{definition}

We note a useful property of the definition given here, which will allow for impredicativity in the model. To do so, we will call a type $A \in \Ty(\Gamma)$ modest if $A_\gamma$ is modest for every $\gamma \in \Gamma$.
\begin{proposition}
	For any context $\Gamma$, if $B \in \Ty(\Gamma.A)$ is modest, then $\Pi(A, B)$ is modest for any $A \in \Ty(\Gamma)$.
\end{proposition}
\begin{proof}
	It suffices to show that $\Pi(A, B)_\gamma$ is modest. Suppose that we have $e \in \bN$ such that $e \Vdash_{\Pi(A, B)_\gamma} f_1, f_2 \in \Pi_{a \in A_\gamma}  B_{(\gamma, a)}$. Then for any $a \in A_\gamma$, there is some $n \in \bN$ such that $n \Vdash_{A_\gamma} a$ and thus $e \cdot n \Vdash_{B_{(\gamma, x)}} f_1(a), f_2(a)$, which shows that $f_1(a) = f_2(a)$ since $B$ is modest, and thus $f_1 = f_2$, which shows that $\Pi(A, B)$ is modest.
\end{proof}
Note that we do not allow general impredicativity in the theory, so we do not make use of this property in the general case when interpreting large dependent products. However, the universal quantifier is interpreted simply as an impredicative product over $\Size$, which is thus validated by the model.

Now we give analogous definitions for the small impredicative product type.
\begin{definition}[Small Product]
	Given a type $A \in \Ty(\Gamma)$ and a small type $B \in \Tm(\Gamma.A, \U)$, we define the small (impredicative) product $\pi(A, B) \in \Tm(\Gamma, \U)$ as a $\Gamma$-indexed set of PERs, where $\pi (A, B)_\gamma$ is the PER defined as follows:
	\begin{align*}
		d &\sim_{\pi (A, B)_\gamma} e \Iff \\ &\forall a \in A_\gamma, \forall m, n \in \bN \. m, n \Vdash_{A_\gamma} a \implies (d \cdot m) \sim_{B_{(\gamma, a)}} (e \cdot n).
	\end{align*}
\end{definition}
	Note that $\pi(A, B)_\gamma$ relates two codes if they are realisers for the same dependent morphism $f : (a \in A_\gamma) \to B_{(\gamma, a)}$ and is thus isomorphic to the set of those dependent morphisms.
    
    The definitions of the term formers also follow from the previous observation, although we can also give explicit definitions.
\begin{definition}[Small term formers]
	Given a term $p : (\gamma \in \Gamma) \to \iota(\pi(A, B)_\gamma)$ and $a \in \Tm(\Gamma, A)$ we define a term $\app^\pi(p, a) :(\gamma \in \Gamma) \to \iota(B(\gamma, a))$ as:
	\[
	\app_{A, B}^\pi(p, a)(\gamma) = \{ d \cdot n \mid d \in p_\gamma \land n \Vdash_{A_\gamma} a_\gamma\}.
	\]
	It is clear from the definition of $\pi(A, B)$ that $\app_{A, B}^\pi(p, a)(\gamma)$ is an equivalence class of the PER $B(\gamma, a)$. 
	Conversely, if we are given a term $t : ((\gamma, a) \in \Gamma.A) \to \iota(B(\gamma, a))$, we define a term $\abs^\pi_{A, B}(t) : (\gamma \in \Gamma) \to \pi(A, B)(\gamma)$ as:
	\begin{align*}
		\abs_{A, B}^\pi&(t)(\gamma) = \\
		&\{ e \mid \forall a \in A_\gamma \. \forall n \in \bN \. n \Vdash_{A_\gamma} a \implies (e \cdot n) \in t_{(\gamma, a)} \}.
	\end{align*}
	Equivalently, $\abs_{A, B}^\pi(t)(\gamma)$ is the set of realisers for the dependent morphism defined by \[a \mapsto t_{(\gamma, a)} : (a : A) \to B(\gamma, a)\]
\end{definition}
These definitions are used to interpret small dependent products. More on the formal interpretation of the dependent product and why the case distinction between small and large products is required can be found in \cite{REUS1999128}.

\subsection{Dependent sum}
The interpretation of the dependent sum requires a similar case distinction as the dependent product. We first give the definition of the large dependent sum type, which is simply a generalisation of context extension, defined as follows.
\begin{definition}[Dependent sum]
	Given a type $A \in \Ty(\Gamma)$ and $B \in \Ty(\Gamma.A)$, the type $\Sigma(A, B) \in \Ty(\Gamma)$ is defined as follows:
	\begin{align*}
		& \Sigma(A, B)_\gamma  = \{ (a, b) \mid a \in A_\gamma \land b \in B_{(\gamma, a)}\}, \\
		& \langle n_1, n_2 \rangle \Vdash_{\Sigma(A, B)_\gamma} (a, b) \Iff n_1\Vdash_{A_\gamma} a \and n_2 \Vdash_{B_{(\gamma, a)}} b.
	\end{align*}
	To interpret the pairing operation, we define a morphism \begin{align*}
		\pair^\Sigma_{A, B}&: \Gamma.A.B \to \Gamma.\Sigma(A, B), \\
		\pair^\Sigma_{A, B}&(\gamma, a, b) = (\gamma, (a, b)).
	\end{align*}
	This morphism is tracked by the pairing operation on the realisers.
	To interpret the elimination rule, we are given a type $P \in \Ty(\Gamma.\Sigma(A, B))$ and a term $p \in \Tm(\Gamma.A.B, P[\pair^\Sigma_{A, B}])$ and define a term $\ind_\Sigma^{p} \in \Tm(\Gamma.\Sigma(A, B), P)$ as follows:
	\begin{align*}
		\ind_\Sigma^{ p}(\gamma, (a, b)) = p_{(\gamma, a, b)}.
	\end{align*}
	If $p$ is tracked by $d$, then this term is tracked by $\Lambda k . \Lambda n \. d \cdot \langle k, \pr_1 n, \pr_2 n \rangle$.
\end{definition}
Note that here we are being somewhat informal in defining the realiser for the elimination rule, since we only assumed a computable bijection $\langle \cdot , \cdot \rangle$ for pairing. However, such a computable bijection may be defined for $n$-tuples as well. 

For the dependent sum type, it is easy to see that if $A$ and $B$ are modest, then $\Sigma(A, B)$ is modest too. We give an analogous definition for the type constructor in $\Per$. Note that, unlike for the $\Pi$-type, we do not require separate term constructors in $\Per$.

\begin{definition}[Small dependent sum]
	Given a small type $A \in \Tm(\Gamma, \U)$ and $B \in \Tm(\Gamma. \El A, \U)$, we define the dependent sum $\sigma (A, B) \in \Tm(\Gamma, \U)$ pointwise as the PER $\sim_{\sigma (A, B)(\gamma)}$ defined as follows:
	\begin{align*}
		\langle m_1, n_1 \rangle \sim_{\sigma(A, B)_\gamma} \langle m_2, n_2 \rangle \Iff m_1 \sim_{A(\gamma)} m_2 \and n_1 \sim_{B(\gamma, [m_1]_{A(\gamma)})} n_2.
	\end{align*}
	Note that this definition is equivalent to $\sigma(A, B) = \E(\Sigma (\El A, \El B))$ and thus we have $\sigma(A, B) \cong \Sigma(\El A, \El B)$. 
\end{definition}

\subsection{Type of sizes}
The impredicativity of our realisability model allows us to interpret the type $\Size$ as a large type, while still having universal quantification over $\Size$ with a small type as codomain live in $\U$.  
Thus, given a context $\Gamma$, we define the semantic type $\Size \in \Ty(\Gamma)$ as $\Size_\gamma = \nabla \omega_1$, where $\omega_1$ is the first uncountable ordinal. Note that the set $\omega_1$ is exactly the set of all countable ordinals. Giving the type a full realisability structure encapsulates the idea that sizes should be treated parametrically.

The operators $\circ \in \Tm(\Gamma, \Size)$ and ${\uparrow}  \in \Tm(\Gamma.\Size, \Size)$ can be interpreted as the constant-zero and successor functions respectively:
\begin{align*}
	\circ(\gamma) &= 0, \\
	\uparrow \!(\gamma, \alpha) &= \alpha + 1.
\end{align*}
Note that since $\Size_\gamma$ carries a full realisability structure, these are tracked by any total realiser. 
We will sometimes treat a term $s \in \Tm(\Gamma, \Size)$ of type $\Size$ as a constant value in $\Size_\gamma$, since we build up such terms from the zero, successor and supremum functions, which behave uniformly over $\Gamma$.

The ordering on sizes $\leq^\U$ is interpreted as a small type in a context with two size variables, i.e. $\leq^\U \,\in \Tm(\Gamma.\Size.\Size[\p_\Size], \U)$. The small type $s_1 \leq^\U s_2$ is either the terminal singleton PER or the initial empty PER, depending on whether $s_1 \leq s_2$ holds with the regular relation on natural numbers.
\begin{align*}
	& \le^U (\gamma, s_1, s_2) = 
	\begin{cases}
		\1 & \text{if } s_1 \le s_2,\\
		\0 & \text{otherwise}.
	\end{cases}
\end{align*}
The actual type $\le \, \in \Ty(\Gamma.\Size.\Size[\p_\Size])$ is then simply the lifting to assemblies: $\le \, = \El (\le^\U)$. We will often write $(s_1 \le s_2)_\gamma$ instead of $\le(\gamma, s_1, s_2)$ and we will blur the distinction between $\le$ and $\le^\U$, since the important property is that $s_1 \le s_2$ is pointwise either the initial or terminal object in the category of PERs or assemblies. Note that this interpretation validates the notion that the ordering on sizes in the theory is a proposition.




\subsection{Bounded quantification and the fixpoint operator}
Recall that the notation $\Gamma, i : \Size \vdash \forall j < i \. A(j)$ is shorthand for $\Gamma, i : \Size \vdash \forall j \. (j < i \to A(j))$. Thus, a term $g$ of type $\forall j < i \. A(j)$ is interpreted as a $\Gamma.\Size$-indexed collection of dependent morphisms, where
\[
g_{(\gamma, \alpha)} : (\beta \in \Size_\gamma) \to ((\beta < \alpha)_\gamma \to A_{(\gamma, \beta)}).
\]

If $(\beta < \alpha)_\gamma $ holds in the interpretation, then $(\beta < \alpha)_\gamma$ is the terminal singleton assembly. Then the set of morphisms of type $(\beta < \alpha)_\gamma \to A_{(\gamma, \beta)}$ is isomorphic to $A_{(\gamma, \beta)}$, allowing us to treat $g_{(\gamma, \alpha)}(\beta)$ as an element of $A_{(\gamma, \beta)}$. 
If $(\beta < \alpha)_\gamma$ does not hold in the interpretation, then $(\beta < \alpha)_\gamma$ is the initial (empty) assembly and there thus exists a unique morphism of type $(\beta < \alpha)_\gamma \to A_{(\gamma, \beta)}$, namely the empty morphism, again given by $g_{(\gamma, \alpha)}(\beta)$. This confirms the syntactic intuition that we may treat $g_{(\gamma, \alpha)}$ as a dependent morphism where $\dom(g_{(\gamma, \alpha)}) = \Size_\gamma \upharpoonright \alpha$. We can thus write the type of $g_{(\gamma, \alpha)}$ as:
\[
g_{(\gamma, \alpha)} : (\beta \in \Size_{<\alpha, \gamma}) \to A_{(\gamma, \beta)}.
\]
Given this intuition, we consider the interpretation of the fixpoint operator $\fix$. Given a type $A \in \Ty(\Gamma. \Size)$, and a dependent morphism
\[
f_{\gamma} : (\alpha : \Size_\gamma) \to ((\beta \in \Size_{<\alpha, \gamma}) \to A_{(\gamma, \beta)}) \to A_{(\gamma, \alpha)},
\]
corresponding to a semantic term of the syntactic type $\forall i . (\forall j < i . A(j) \to A(i))$,
we need to define a term $\fix f : (\gamma \in \Gamma) \to \Pi (\Size, A)_\gamma$. Thus, we require that $(\fix f)_\gamma$ is a dependent morphism from $\Size$ to $A$, i.e. 
\[
(\fix f)_\gamma : (\alpha \in \Size_\gamma) \to A_{(\gamma, \alpha)}.
\] 
We define this morphism by well-founded induction on $\Size_\gamma$, i.e. the ordinal $\omega_1$ and thus a well-founded ordering:
\begin{align*}
	(\fix f)_\gamma(\alpha) &= f_\gamma(\alpha, \lambda \beta < \alpha \. (\fix f)_\gamma(\beta)).
\end{align*}
In the case where $\alpha = 0$, there is only the unique empty morphism with type $(\beta \in \Size_{<0, \gamma}) \to A_{(\gamma, \beta)}$, since the set $\Size_{<0, \gamma}$ is empty. 
Note that this definition mirrors exactly the $\beta$-rule for $\fix\!$. The uniqueness property for $\fix f$ also follows by well-founded induction over $\Size_\gamma$. 

For $\fix$ to be a term, it must be tracked by some realiser. Note that, by our definition of $\fix$, if we are given some $f$ realised by $f_r \in \bN$ and $\gamma$ realised by $\gamma_r \in \bN$, we require that a realiser $F$ for $\fix$ satisfying the following for every $n \in \bN$, where we use as convention that the partial operation associates to the left so that we may omit parentheses:
\[
F \cdot f_r \cdot \gamma_r \cdot n = f_r \cdot \gamma_r \cdot n \cdot (\Lambda m \. F \cdot f_r \cdot \gamma_r \cdot m).
\]
In order to define such a realiser, we make use of the fact that a fixpoint operator exists in every partial combinatory algebra:
\begin{proposition}[Fixpoint operator in pca]
    For every pca $(\cA, \cdot)$, there exists a $\fix\! \in \cA$ such that 
    \[
    \fix f \downarrow \text{ and } 
    \fix f \, a \simeq f \, (\fix f) \, a.
    \]
\end{proposition}
A proof of this can be found in \cite{oosten}.
Importantly, it allows us to define a realiser for the fixpoint operator. First, we define a code $\varphi \in \bN$ by
\[
\varphi \coloneq  \Lambda f. \Lambda e . \Lambda k . \Lambda n \. e \cdot k \cdot n \cdot (\Lambda m \. f \cdot e \cdot k \cdot m).
\]
This gives us the following:
\begin{align*}
(\fix \varphi) \cdot f_r \cdot \gamma_r \cdot n &\simeq \, \varphi \cdot (\fix \varphi) \cdot f_r \cdot \gamma_r \cdot n \\
&\simeq f_r \cdot \gamma_r \cdot n \cdot (\Lambda m . (\fix \varphi) \cdot f_r \cdot \gamma_r \cdot m).
\end{align*}
This shows that $\fix \, \varphi$ is a realiser for our fixpoint operator.



\subsection{Existential type}
To interpret the existential type $\exists i . A(i)$, we quotient the corresponding $\Sigma$-type in order for it to live inside the universe $\U$, i.e. canonically construct a partial equivalence relation from the assembly $\Sigma(\Size, A)$. To do so, we make use of the following lemma.

\begin{lemma}
The inclusion functor $I : \Mod \to \Asm$ has a left adjoint $\M : \Asm \to \Mod$.
\end{lemma}
\begin{proof}
Let $A \in \Asm$ be an assembly, we define $\M(A) = ( A {/ \!\sim_\M}, \Vdash_{\M(A)})$ where $\sim_\M$ is the least equivalence relation such that the following holds:
\[
\exists n \in \bN \. n \Vdash_A a_1, a_2 \implies a_1 \sim_\M a_2.
\]
The realisability relation is defined as
\[
n \Vdash_{\M(A)} [ a ]_{\sim_\M} \Iff \exists a' \in [ a ]_{\sim_\M} (n \Vdash_A a').
\]
It is clear that $\M(A)$ is indeed a modest set. We define a map $\eta_A : A \to \M(A)$ as $a \mapsto [a]_{\sim_\M}$, which is realised by $\Lambda x . x$ and is thus a morphism of assemblies. To show that $\M$ extends to a left adjoint of $I$, let $X \in \Mod$ be a modest set and suppose there exists a morphism $g : A \to X$. Then $g$ uniquely factors as $f \circ \eta_A$, where $f : \M(A) \to X$ is defined as $f([a]_{\sim_\M}) = g(a)$. 

It remains to show that $f$ is well-defined. Suppose that $a_1 \sim_\M a_2$; if there exists $n \in \bN$ such that $n \Vdash_A a_1, a_2$, then since $g$ is tracked by some $e \in \bN$, this means that $e \cdot n \Vdash_X g(a_1), g(a_2)$. Thus, since $X$ is a modest set, $g(a_1) = g(a_2)$ and so $f$ is well-defined.

Else, there exist $n_1, ..., n_k \in \bN$ and $b_1, ..., b_{k-1} \in A$ such that $n_1 \Vdash_A a_1, b_1$, $n_2 \Vdash_A b_1, b_2$, ..., and $n_k \Vdash_A b_{k-1}, a_2$. Then by similar reasoning via transitivity, we will get that $g(a_1) = g(b_1) = \dots = g(b_{k-1}) = g(a_2)$.
\end{proof}

Given this construction, we can truncate any assembly $A \in \Asm$ to $\M(A)$ in order (almost) fit in the universe $\U$. However, the universe $\U$ consists of PERs rather than modest sets. Fortunately, the categories $\Per$ and $\Mod$ are equivalent, allowing us define the \emph{$U$-truncation} $\|A\| \in \Tm(\Gamma, \U)$ of a type $A \in \Ty(\Gamma)$ as 
\[\Tr{A}_\gamma = \M(A_\gamma).\]

The $U$-truncation operation should also operate on the level of terms, that is, given a term $a \in \Tm(\Gamma, A)$, there should be a term $\tr{a}$ in the $U$-truncated type. Thus, we define a context morphism 
\begin{align*}
	&\tr{\cdot} : \Gamma.A \to \Gamma.\El (\Tr A), \\
	&\tr{(\gamma, a)} = (\gamma, [a]_{\sim_\M}),
\end{align*}
which is realised by the identity $\Lambda x . x$.
Moreover, the existential type comes with an elimination rule for small types. Thus, a truncated type $\Tr A$ in the model should come with such a rule as well. That is, given a small type $P \in \Tm(\Gamma.\El \Tr A, \U)$ and a term $t \in \Tm(\Gamma.A, (\El P)[\tr \cdot ])$, we define the following term
\begin{align*}
	&\elim_t \in \Tm (\Gamma.\El (\Tr A), \El P), \\
	&\elim_t (\gamma, [a]_{\sim_\M}) = t(\gamma, a).
\end{align*}
To show that $\elim_t$ is well defined, we need to show that $a_1 \sim_M a_2$ implies $t(\gamma, a_1) = t(\gamma, a_2)$. Since $a_1 \sim_M a_2$, we can assume that there exists some $k \in \bN$ such that $k \Vdash_A a_1, a_2$. Moreover, $t$ is tracked by some $e \in \bN$, meaning that the following holds:
\[
\langle m, n \rangle \Vdash_{\Gamma.A} (\gamma, a) \implies e \cdot \langle m, n \rangle \Vdash_{P(\gamma, [a]_{\sim_\M})} t(\gamma, a).
\]
Then for any $\gamma \in \Gamma$, which must realised by some $m \Vdash_\Gamma \gamma$, we have that $\langle m, k \rangle \Vdash_{\Gamma.A} (\gamma, a_1), (\gamma, a_2)$ and thus:
\[
e \cdot \langle m, k \rangle \Vdash_{P(\gamma, [a_1]_{\sim_M})} t(\gamma, a_1), t(\gamma, a_2), 
\]
since $[a_1]_{\sim_M} = [a_2]_{\sim_M}$. Since $P$ is a small type, $P(\gamma, [a_1]_{\sim_M})$ is a modest set when viewed as an assembly. This implies that $t(\gamma, a_1) = t(\gamma, a_2) $ since any realiser codes at most one object. 

Now we are ready to interpret the existential type by truncating the $\Sigma$-type.
\begin{definition}[Existential Type]
	Given a small type $A \in \Tm(\Gamma.\Size, \U)$, the small existential type $\exists i . A(i) \in \Tm(\Gamma, \U)$ is defined as 
	\[\exists i . A(i) \coloneq  \Tr{\Sigma(\Size, \El A)}.\]
	
	As with dependent sums, the pairing operation is interpreted with a morphism $\pair_\exists : \Gamma.\Size.A \to \Gamma. \El (\exists i.A(i))$ defined as 
	\[
	\pair_\exists(\gamma, s, a) = \tr{\pair_\Sigma(\gamma, s, a)},
	\]
	where $\pair_\Sigma(\gamma, s, a) \in \Gamma.\Sigma(\Size, A)$.

	Finally, the elimination rule for existential types is interpreted as follows. Given the following terms:
	\begin{align*}
		P &\in \Tm(\Gamma. \El (\exists i.A(i)), \U),\\
		p &\in \Tm(\Gamma.\Size.A, (\El P)\{\pair_\exists\}),
	\end{align*}
    we define the term $\ind_\exists^{p} \in \Tm(\Gamma.(\exists i. A(i)), \El P)$ as follows:
	\[
		\ind_\exists^{p} = \elim_{\ind_{\Sigma}^{P}} \in \Tm(\Gamma.(\exists i. A(i)), \El P).
	\]
	where $ \elim_{\ind_{\Sigma}^{P}}(\gamma, [(s, a)]_{\sim_\M}) = \ind_{\Sigma}^{P}(\gamma, (s, a))$ is well-defined since $\elim$ is well-defined.

\end{definition}

Now that the formal definition is given, we can make an observation about the structure of the existential type. Given a small type $A \in \Tm(\Gamma, \U)$, note that for the modest set $\M (\Sigma(\Size, \El A)_\gamma)$, the equivalence relation $\sim_\M$ is defined exactly by
\begin{align*}
	(s_1, a_1) \sim_\M (s_2, a_2) &\Iff \exists \langle m, n \rangle \in \bN \. \langle m, n \rangle \Vdash_{\Sigma(\Size, \El A)_\gamma} (s_1, a_1), (s_2, a_2) \\
	&\Iff \exists n \in \bN \. n \Vdash_{\El A_{(\gamma, s_1)}} a_1 \and n \Vdash_{\El A_{(\gamma, s_2)}} a_2
\end{align*} 
Thus, the intuition is that two elements $[(s_1, a_1)]_{\sim_\M}$ and $[(s_2, a_2)]_{\sim_\M}$ of $(\exists i . A(i))_\gamma$ are equal if $a_1$ and $a_2$ share a realiser.